\documentclass[protokoll]{dvd}

\KOMAoptions{
    headwidth = 18cm,
    footwidth = 18cm,
}
\usepackage{graphicx}
\begin{document}

\title{Styrelsemöte 5}
\subtitle{2023}
\author{Styrelsen}
\date{2023-11-26}


\textbf{Datum:} \csname @date\endcsname\\
\textbf{Tid:} 16:10\\
\textbf{Plats:} Monaden\\
\textbf{Styrelsemedlemmar:}
\begin{närvarande_förtroendevalda}
    \förtroendevald{Ordförande}{Samuel Hammersberg}{Ja}
    \förtroendevald{Vice ordförande}{Vakant}{Nej}
    \förtroendevald{Kassör}{Lukas Gartman}{Ja}
    \förtroendevald{SAMO}{Tekla Siesjö}{Nej}
    \förtroendevald{Sekreterare}{Vakant}{Nej}
\end{närvarande_förtroendevalda}

\section{Öppnande av möte}

Mötet öppnades av Samuel Hammersberg kl 16:10

\section{Runda bordet}

Runda bordet innebär att varje person berättar hur de känner sig.
Man kan till exempel berätta att man är stressad på grund av en inlämning, irriterad på sin granne, eller bara väldigt glad därför att man ligger i fas med plugget.

\section{Formalia}

\subsection{Styrelsens beslutbarhet}

\blockquote[7 kap. 5 \S~första stycket i stadgan][]{%
    Styrelsen är endast beslutsmässig då samtliga styrelsemedlemmar har fått kallelsen till styrelsemötet och minst hälften av styrelsemedlemmarna är närvarande.
    Ordförande eller vice ordförande måste vara närvarande när beslut tas.
}

\subsubsection*{Förslag}

\begin{attsatser}
    \item Styrelsen har uppnått kraven i 7 kap. 5 § första stycket i stadgan och är därmed beslutbar.
\end{attsatser}
\subsubsection*{Beslut}
\begin{attsatser}
    \item förslaget till beslut bifalles
\end{attsatser}


\subsection{Fastställande av mötesschema}

För att styrelsen ska kunna fatta ett styrelsebeslut eller protokollföra en diskussion behöver punkten i mötesschemat där styrelsen ska fatta beslut vara inlagd eller föras in i mötesschemat senast vid den här punkten.

\subsubsection*{Förslag}

\begin{attsatser}
    \item mötesschemat fastställs utan några förändringar.
\end{attsatser}
\subsubsection*{Beslut}
\begin{attsatser}
    \item förslaget till beslut bifalles
\end{attsatser}

\subsection{Val av mötessekreterare}
Eftersom att styrelsen för nuvarande inte har en sekreterare tar Lukas Gartman på sig jobbet.
\subsubsection*{Förslag}
\begin{attsatser}
    \item Lukas Gartman väljs till mötessekreterare
\end{attsatser}
\subsubsection*{Beslut}
\begin{attsatser}
    \item förslaget till beslut bifalles
\end{attsatser}

\subsection{Val av protokolljusterare}

Protokolljusterare har till uppgift att kontrollera att protokollet i slutändan reflekterar de faktiska besluten och diskussionerna som fördes under mötet.
Utöver protokolljusteraren så ska mötesordförande och mötessekreteraren signera protokollet.
Vid styrelsemöten ska det endast vara en justerare.
Mötesordförande och mötessekreteraren kan inte vara justerare.

\subsubsection*{Förslag}
\begin{attsatser}
    \item Albin Otterhäll väljs till protokolljusterare
\end{attsatser}
\subsubsection*{Beslut}
\begin{attsatser}
    \item förslaget till beslut bifalles
\end{attsatser}

\section{Rapport}
\subsection{Ordförande}
Samuel Hammersberg har fixat access för att låsa upp dörren för Tim Persson, då han ska kunna låsa upp dörren under mottagningen.
Utöver detta har han varit med och hjälpt

\subsection{Kassör}
Lukas Gartman har kontaktat instutionen för att få reda på status på när mottagningspengarna kommer in.
Tydligen har det varit strul men pengarna ska komma in den 22:e eller den 23:e augusti.

\newpage

\section{Rapporter}
\subsection*{Ordförande}
Samuel Hammersberg har suttit i tre möten om flytten av Lindholmen campuset.
Under dessa möten har det för det mesta diskuteras adminstrativa ämnen, såsom 
hur kommunikation ska fungera, men i alla fall minstånde så är Monadens 
kårtillförighet försäkrats och den kommer fortsätta ligga under DVD.

Han har även varit två Instutionsråd då det diskuterades bland annat hur man ska
kunna anställa mer amanueser. Samuel har hjälpt till med detta och skickat ut reklam 
för tjänsterna.

Samuel har även pratat med Gustav Dalemo, Tim Persson, och Josefin Kokkinakis
om deras intresse av att vara delaktiga i styrelsen!

\subsection*{Kassör}

\section{Beslutspunkter}

\subsection{Fastställande av kommande stämmas mötesschema}
Då styrelsen ska hålla i en stämma den 30de November 2023, så ska styrelsen nu
bestämma de mötespunkter de vill ta upp på stämman.

Nuvarande styrelsen kommer skriva ihop en rapport så att de kan presentera 
deras arbete sedan senaste stämman.

Under året har styrelsen märkt att Tim Persson, Gustav Dalemo och Josefin Kokkinakis
hade passat som styrelsemedlemmar. Gustav Dalemo har visat intresse för att vara 
sekreterare, Tim Persson har visat intresse för vice-ordförande och Josefin Kokkinakis 
för rolen som SAMO. Styrelsen kommer även att ha ett möte med dom innan stämman för att
kolla av intresset.

Eftersom att Götas IT-Sektion har bestämt på en budget av 1000kr per förening,
så känner styrelsen att hur detta belopp ska spenderas bör diskuteras under stämman.

På grund av att styrelsen ej riktigt jobbat på väggarbetet, bör även detta tas upp.

Styrelsen hade även velat se till att vi kan få våran logga målad på Olgas trappor.

\subsubsection*{Förslag till beslut:}
\begin{attsatser}
    \item lägga till punkter för rapporter från styrelsen, nya styrelsemedlemmar,
          diskussion över hur våra äskningar gentemot Göta ska fungera under året,
          väggarbetet, måla divisionens inofficella logga vid Olgas trappor.
\end{attsatser}

\subsubsection*{Beslut}
\begin{attsatser}
    \item förslaget till beslut bifalles
\end{attsatser}


\section{Diskussion}

\newpage
\section{Avslutande av möte}

\subsection{Nästa möte}
Styrelsen har preliminärt bokat in 31:a 10:00 för nästa möte.

\subsection{Mötets avslutande}
Mötet avslutades kl. 15:00

\styrelsesignaturer

\end{document}
